\section{Novelty and Contributions}
\label{sec:novelty}

This section states precisely what is new in UCRA relative to the
prior art surveyed in Section~\ref{sec:related}, and what is
deliberately \emph{not} claimed. We organize the novelty into five
claims, each paired with the evidence that supports it in
Section~\ref{sec:results}.

\subsection{Claim 1: A single-equation, operator-auditable
uncertainty-to-reservation transform}

Prior forecast-driven provisioning either buries the uncertainty
inside a learned policy (RL agents, end-to-end neural controllers) or
uses fixed engineering factors disconnected from the forecast's own
uncertainty estimate. UCRA's transform $\Phi$ (Eq.~\eqref{eq:phi}) is,
to our knowledge, the first slicing reservation rule that (i) consumes
the forecast distribution directly (a quantile and its dispersion),
(ii) modulates the buffer by an online risk indicator computed from
reservation feedback, and (iii) exposes exactly one interpretable
parameter $\kappa$ whose effect traces a measurable risk--utilization
frontier (Fig.~\ref{fig:kappa}). The transform is auditable in the
strong sense of Rudin~\cite{rudin2019stop}: given any reservation
decision, the operator can decompose it into forecast, uncertainty,
risk, and policy terms in seconds.

\subsection{Claim 2: Closed-loop coupling with churn suppression}

Stage~4's update rule couples the reservation target to realized
demand through EWMA smoothing plus \emph{release hysteresis}: slack is
released only after sustained non-use. On real 15-minute RAN data this
combination yields reservations that track demand (10.7\% tracking
error) while never violating, and it suppresses the oscillation that
naive target-chasing produces on bursty traffic. The combination of
probabilistic targets, hysteresis-bounded actuation, and per-slot
feedback into both the transform and the drift monitor is, to our
knowledge, not present as a unit in prior slicing resource managers.

\subsection{Claim 3: A formally causal protocol for online-learning
evaluation}

The self-evolution stage is architecturally conventional (drift
trigger + reservoir replay + fine-tuning); the novelty is the
evaluation contract. We formalize \emph{causal label availability} for
multi-step forecasters: a window observed at time $t$ with horizon
$H$ may enter the training pool no earlier than $t+H-1$
(Proposition~\ref{prop:causal}), and a model update at time $t$ may
only alter forecasts for slots later than $t$. We are not aware of
prior slicing-reservation work that states, enforces, and unit-tests
this contract. The practical significance is demonstrated
empirically: the leaky variant of our own pipeline inflates the
apparent benefit of self-evolution, and the corrected protocol still
supports the qualitative claim (halved post-surge violations) with
defensible numbers. We offer the protocol as a reusable standard for
evaluating self-adaptive network controllers, complementary to the
general leakage taxonomy of Kapoor and Narayanan~\cite{kapoor2023leakage}.

\subsection{Claim 4: Trigger discipline as a first-class result}

Self-evolving systems are usually evaluated only on scenarios where
adaptation helps. UCRA is evaluated symmetrically. On the clean RAN
window --- three weeks of live traffic --- the drift triggers fire
\emph{zero} times in every seed while the system holds 0.0\%
violations: the insurance policy is armed but never cashed, and never
fires spuriously. On the surged window the same triggers fire
reliably and adapt. This symmetric evidence (silence when unneeded,
action when needed) is, we argue, the correct standard for
evaluating trigger-based adaptivity, and it is rarely reported.

\subsection{Claim 5: Verified reproducibility as a methodological
contribution}

Every headline number in this paper is regenerated by a scripted
pipeline from pinned public datasets, with configuration-controlled
seeds, a three-seed verification driver, causality regression tests,
and an explicit leakage audit (Section~\ref{sec:leakage-audit}).
During preparation, the audit invalidated two of our own earlier
results (the replay leakage above, and an outcome-informed risk term
in the CTTC evaluation), which we document and correct in the final
version. We consider the audit methodology --- treat your own
information flows as adversarially reviewable --- a contribution in
its own right for the ML-for-networking community, where evaluation
subtleties are easy to miss and expensive to discover post
publication.

\subsection{What we do not claim}

UCRA does not claim a novel forecaster: the quantile LSTM is a
deliberately standard, well-understood architecture
(Section~\ref{sec:stage1}), and forecast accuracy is treated as a
solved-enough subproblem. We do not claim optimality of $\kappa$; the
frontier (Fig.~\ref{fig:kappa}) is presented so each operator can
choose. We do not claim the CTTC risk-modulation term is decisive ---
on that dataset the honest finding is that it is not
(Section~\ref{sec:results-cttc}), and we say so. Finally, we do not
claim generality beyond the evaluated regimes; threats to validity
are cataloged in Section~\ref{sec:discussion}.
